\documentclass[11pt]{article}

% Change "review" to "final" to generate the final (sometimes called camera-ready) version.
% Change to "preprint" to generate a non-anonymous version with page numbers.
\usepackage{acl}

% Standard package includes
\usepackage{times}
\usepackage{lipsum}
\usepackage{latexsym}

% For proper rendering and hyphenation of words containing Latin characters (including in bib files)
\usepackage[T1]{fontenc}
% For Vietnamese characters
% \usepackage[T5]{fontenc}
% See https://www.latex-project.org/help/documentation/encguide.pdf for other character sets

% This assumes your files are encoded as UTF8
\usepackage[utf8]{inputenc}

% This is not strictly necessary, and may be commented out,
% but it will improve the layout of the manuscript,
% and will typically save some space.
\usepackage{microtype}

% This is also not strictly necessary, and may be commented out.
% However, it will improve the aesthetics of text in
% the typewriter font.
\usepackage{inconsolata}

%Including images in your LaTeX document requires adding
%additional package(s)
\usepackage{graphicx}

% If the title and author information does not fit in the area allocated, uncomment the following
%
%\setlength\titlebox{<dim>}
%
% and set <dim> to something 5cm or larger.

\title{Scaling Multi-Turn RAG: High-Performance Parallelized Pipeline for the MTRAG Benchmark}

\author{Murat Çelik \quad Nejla Dinçer \quad Can Ersoy \\
        \textbf{Mert Toprak} \quad \textbf{Barış Tan Ünal} \\
        Politecnico di Torino \\
        Corso Duca degli Abruzzi, 24 10129 Torino, Italy \\
        \texttt{\{s338212, s343290, s336887, s336966, s338210\}@studenti.polito.it}}

\begin{document}
\maketitle
\begin{abstract}
Recently, Retrieval-Augmented Generation (RAG) has become a significant task in Large Language Models (LLMs). In multi-turn Retrieval-Augmented Generation (RAG), a good system must overcome the challenges of maintaining context as the dialogue turns progress and manage the issue of generating answers depending on conversation history. Our project presents a high-performance, parallelized Multi Turn RAG pipeline designed for systematic evaluation of the MTRAG benchmark across its three subtasks: Retrieval (Task A), Generation (Task B), and End-to-End RAG (Task C). Our methodology utilizes a Streamlit framework that allows users to embed diverse corpora with varying vector spaces and embedding models, facilitating configuration for each task regarding the task’s nature. Some key experiments are focused on the performances of different vector databases, embedding models, the necessity of LLM-based query rewriting for non-standalone questions, utilization of different rerankers and scale and performance of the selected LLM for the answer generation. We conclude that a configuration utilizing query rewriting with reranking delivers the better results comparing to their absence. You can find our GitHub repository \href{https://github.com/merttoprak1/MTRAGEval-Evaluating-Multi-Turn-RAG-Conversations}{here}.
\end{abstract}

\section{Introduction}
Large Language Models (LLMs) have revolutionized information-seeking tasks, yet they often struggle with context preservation, hallucination and generating answers given a corpora. Retrieval-Augmented Generation (RAG) addresses this by basing the model responses on relevant external passages. While single-turn RAG is becoming a immerse field, Multi-Turn RAG presents a more complex challenge where the system must generate responses to a question within the context of an ongoing conversation. Unlike single-turn, multi-turn interactions require the LLM to track conversational history and evolve its retrieval strategy as the dialogue progresses. One of the main causes of the challenge is because of the non-standalone nature of human dialogue. In later turns of a conversation, users frequently use shortcuts, such as referencing (e.g., "its address?" or "the previous one"), which make the current question ambiguous without the preceding context. Moreover, the system requires active retrieval strategies, where the it must decide whether a question is answerable based on the available corpus it is provided with or if it should decline to answer to avoid hallucinations.


The goal of our project is to develop a modular, parallelized RAG pipeline specifically optimized for the MTRAG benchmark. Our system is designed to efficiently process 842 distinct tasks across four diverse domains: CLAPNQ, FIQA, Govt, Cloud. By leveraging a Streamlit-based framework, we aim to maximize performance metrics such as nDCG and Recall for retrieval and RBalg, RBllm and RLf for generation. This architecture allows for a systematic exploration of the performance comparison for the following based on the aforementioned metrics: Varying the utilization of different embedding models, vector databases, existence of query rewriting, existence and utilization of reranking with different top-k values and use of different LLMs for generation. 
The MTRAG Benchmark contains three subtasks: 

\textbf{Task A (Retrieval): } Evaluating the system's ability to retrieve gold-standard passages for the final turn of a conversation.

\textbf{Task B (Generation): } Assessing the generator's ability to produce a faithful answer when provided with perfect reference passages.

\textbf{Task C (End-to-End RAG): } Evaluating the complete pipeline, where the system must perform its own retrieval before generating a final response for the user’s last question.

Our work focuses on providing an system such that the user can dynamically try different configuration settings so that they can receive better metric results. We have also provided an evaluation section such that it is easily possible to visualize the outcomes of their prediction results for each task. We have tested and evaluated the MTRAG Benchmark dataset with various configurations, which will be analyzed in depth in the following sections.

\section{Related Work}
\subsection{Multi-Turn RAG Evaluation}
Recent advancements in Large Language Models (LLMs) have necessitated robust evaluation frameworks. While benchmarks like MT-Bench \citet{zheng2023judging} evaluate multi-turn conversation capabilities using LLM-as-a-judge, they often lack a specific focus on Retrieval-Augmented Generation (RAG) challenges such as domain-specific retrieval and unanswerable questions. The MTRAG Benchmark \citet{katsis2025mtrag} addresses this gap by stressing the system's ability to handle unanswerable queries. Unlike static, single-turn datasets like Natural Questions or MS-MARCO where an answer is always guaranteed, MTRAG requires the model to discern when retrieved contexts are insufficient and explicitly abstain from answering, thereby testing its robustness against hallucination.


\subsection{Retrieval Strategies}
The retrieval component is the backbone of any RAG system. Traditional Sparse Retrieval methods, such as BM25 \citet{Robertson2009}, rely on exact keyword matching. While efficient, they suffer from the lexical gap problem, failing to caption semantic synonyms. Conversely, Dense Retrieval \citet{karpukhin2020dense} employs dual-encoder architectures to map queries and documents into a shared latent space, capturing semantic similarity. Modern state-of-the-art pipelines often employ a hybrid approach or a two-stage "Retrieve-and-Rerank" architecture. In this setup, a lightweight retriever fetches a broad set of candidates, which are then re-scored by a Cross-Encoder \citet{nogueira2019} to maximize precision at the cost of higher latency.

\subsection{Query Rewriting in Conversational AI}
A unique challenge in multi-turn RAG is the ambiguity of user queries due to coreferences (e.g., "how much is it?" referring to a previously mentioned product). Direct retrieval using such raw queries often leads to poor performance. Query Rewriting (QR) addresses this by generating a standalone, de-contextualized query using the conversation history \citet{vakulenko2021qr}. This rewritten query effectively bridges the gap between conversational nuances and the need for explicit search terms in vector databases.

\section{Dataset}
The MTRAG benchmark aggregates data from four diverse domains to ensure robust evaluation across different linguistic styles and knowledge types: \\
    \textbf{1. ClapNQ:} Adapted from the Natural Questions dataset, this domain focuses on open-domain factoid questions based on Wikipedia articles. It serves as a general-knowledge baseline. \\
    \textbf{2. Cloud:} Comprising technical documentation from major cloud service providers. This domain tests the system's ability to handle jargon-heavy, structured technical content. \\
    \textbf{3. FiQA:} Financial dataset from forums and reports. It requires numerical reasoning and specialized terminology. \\
    \textbf{4. Govt:} Sourced from various government information portals. This domain involves parsing bureaucratic language and complex policy details.


\section{Methodology}
 Our system adopts a modular, component-based architecture designed for flexibility and high performance. The core pipeline consists of five sequential stages: \\
    \textbf{1. User Input Processing:} The system accepts a raw user query along with the conversation history. \\
    \textbf{2. Query Rewriting:} An LLM-based rewriter transforms the context-dependent query into a standalone search query. \\
    \textbf{3. Embedding Generation:} The rewritten query is converted into a dense vector embedding. \\
    \textbf{4. Vector Retrieval:} A high-speed vector database (Qdrant or Faiss) retrieves the top-$K$ most similar document chunks. \\
    \textbf{5. Reranking:} A cross-encoder model re-scores the retrieved candidates to improve precision, filtering down to the top-$N$ context passages. \\
    \textbf{6. Generation:} The final context is fed into the LLM generator to produce the answer.
    
    
    \subsection{Task A: Retrieval Only}
    Task A isolates the retrieval subsystem. Given a user query and conversation history, the system must retrieve the specific passages that contain the answer.
    \textbf{Query Rewriting (QR):} A critical challenge in multi-turn conversations is detecting the intent of elliptical or coreferential queries (e.g., "Why did he do that?"). We employ a generative Query Rewriter that leverages the conversation history to resolve these ambiguities. The rewriter is prompted to produce a fully specified, standalone question that captures the user's original intent without requiring external context, which is essential for effective vector similarity search.
    
\textbf{Vector Search Infrastructure:} We evaluated two distinct vector search backends:

\textbf{Faiss:} A library for efficient similarity search and clustering of dense vectors. It provides a lightweight, in-memory solution suitable for rapid prototyping but lacks advanced filtering features.

\vspace

\textbf{Qdrant:} A production-grade vector database that supports HNSW (Hierarchical Navigable Small World) indexing for approximate nearest neighbor search. Crucially, Qdrant allows for payload-based filtering, enabling efficient retrieval within specific subsets of the corpus.
    
    \textbf{Two-Stage Retrieval \& Reranking:}
    To balance recall and precision, we implement a two-stage retrieval process. In the first stage, the vector database retrieves a broad set of candidates (e.g., Top-$K=20$) based on cosine similarity. In the second stage, a Cross-Encoder model (such as BGE-Reranker or MS-MARCO) performs a computationally intensive re-scoring of these candidates. This reranking step is vital for filtering out "hard negatives"—documents that are semantically similar but factually irrelevant—thereby presenting the generator with only the highest-quality context (Top-$N=10$).

\subsection{Task B: Generation with Reference Passages}
For Task B generation, our main goal is to generate relevant answers to the last question in the given set of tasks. 
Task is defined as a conversation turn containing all previous turns together with the last user question (e.g., the task created for turn $k$ includes all user and agent questions/responses for the first $k−1$ turns plus the user question for turn $k$).

To start generation process, the system converts the MTRAG input format into LangChain-compatible message objects. Then,
all turns except the last one are converted to HumanMessage(user) or AIMessage(agent) objects and stored as conversation history. Subsequently, the last user turn is extracted as the current question to be answered.
\vspace

\textbf{Context Management:}
Retrieved passages are formatted into a structured context  string. Each passage is numbered sequentially and presented in order of the relevance score(descending order).
Afterwards, the system constructs a comprehensive prompt concatenating three key components:
\textbf{System Prompt:} Defines the assistant's(LLM model) role and instructions.
\textbf{Conversation History:} All previous turns (HumanMessage/AIMessage pairs).
\textbf{Current Question:} The latest user query.

After assembling all the parts for the prompt, it is then sent to the selected LLM model to generate an answer for the user's question.

\textbf{Model Selection:} We used a variety of LLM models to compare which give more reasonable answers. We selected models from different companies and with different sizes to see variations in the generation results between bigger/smaller models and different vendors. The models we used for generation are namely; \texttt{gpt-oss-20b} from openai, \texttt{gemma-3-12b} from google,
 \texttt{deepseek-r1-0528-qwen3-8b} from deepseek, \texttt{ministral-3-14b-reasoning} and \texttt{/ministral-3-3b} from mistralai.

\subsection{Task C: Generation with Retrieved Passages}

Task C combines retrieval and generation in a single pipeline. Given a conversation and a document corpus, the system first retrieves relevant passages and then generates an answer based only on those passages.

\textbf{Query Rewriting:} The last user question may contain unclear references or missing details. We use a Query Rewriter that looks at the full conversation and rewrites the question into a clear, standalone form. This helps improve retrieval accuracy.

\textbf{Passage Retrieval:} The rewritten query is converted into a vector and searched against the document corpus using a vector database (Qdrant). We retrieve the Top-$K=20$ most similar passages to ensure high recall.

\textbf{Reranking and Selection:} A cross-encoder reranker rescales the relevance of the retrieved passages. Based on these scores, we select the Top-$N$ passages. We return at most $N=5$ or $N=10$ passages, which provides enough context while reducing noise.

\textbf{Context Building:} The selected passages are ordered by relevance and combined into a single context block. This context is given to the generation model.

\textbf{Answer Generation:} The final prompt is formed by combining system instructions, conversation history, retrieved context, and the current user question, and is then passed to the LLM to generate the final answer.



This prompt is sent to the LLM, which generates the final answer using only the retrieved passages, ensuring factual and faithful responses.

\subsection{Evaluation Framework}
\begin{itemize}
    \item \textbf{Retrieval Metrics:} Define Recall@K and nDCG@K used for Tasks A and C.
    \item \textbf{Generation Metrics:} Define the "Judge" system (LLM-as-a-Judge):
    \begin{itemize}
        \item \textbf{Faithfulness ($RL_F$):} Measuring if the answer is derived solely from the provided context.
        \item \textbf{IDK Accuracy:} Measuring the model's ability to correctly identify unanswerable questions.
    \end{itemize}
\end{itemize}

\section{Experimental Results}



\subsection{Experimental Setup}
\begin{itemize}
    \item List specific hardware (GPUs) and configurations used in the experiments.
    \item \textbf{Embedding Models:} \texttt{all-minilm:l6-v2}, \texttt{nomic-embed-text}, \texttt{jina-embeddings-v3}.
    \item \textbf{Rerankers:} \texttt{BGE-v2-M3}, \texttt{Qwen3-Reranker}, \texttt{ms-marco}.
    \item \textbf{Parameters:} Top-K retrieval counts (20 vs 30).
\end{itemize}

\subsection{Task A: Retrieval Performance}
\label{subsec:task_a_results}

We evaluated the retrieval pipeline based on Normalized Discounted Cumulative Gain (nDCG@k) and Recall (R@k) at varying depths ($k=3, 5$). Our primary objective was to maximize the relevance of context provided to the generation phase while maintaining low latency. We established a \textbf{Baseline} configuration using Qdrant, \texttt{all-minilm-l6-v2} embeddings, and no reranking, which yielded an nDCG@3 of 0.249.

\subsubsection{Query Rewriting Strategies}
Given the conversational nature of MTRAG, where queries often depend on previous turns (e.g., ``What about in 2024?''), direct retrieval often fails. We compared two models for Query Rewriting (QR): \texttt{google/gemma-3-12b} and \texttt{openai/gpt-oss-20b}.
While enabling QR with the smaller Gemma model improved the baseline nDCG@3 to 0.268 (+7.6\%), the larger GPT-20b model achieved 0.293 (+17.6\%). This substantial gap indicates that the reasoning capabilities required to disambiguate complex multi-turn references scale significantly with model size. Consequently, we utilized GPT-20b for all subsequent experiments to isolate downstream architectural effects.

\subsubsection{Impact of Reranking Architectures}
To mitigate the precision limitations of dense vector retrieval, we introduced a second-stage reranking step. We experimented with cross-encoder architectures to re-score the top-$k$ documents retrieved by the vector store.

\begin{table}[h]
\centering
\small
\begin{tabular}{lcccc}
\hline
\textbf{Reranker Model} & \textbf{Depth} & \textbf{nDCG@3} & \textbf{R@3} \\
\hline
\textit{None (Baseline)} & - & 0.293 & 0.275 \\
MS-Marco-MiniLM & 10 & 0.293 & 0.275 \\
Mxbai-xsmall-v1 & 10 & 0.317 & 0.300 \\
BGE-v2-M3 & 10 & 0.329 & 0.311 \\
\textbf{BGE-v2-M3} & \textbf{5} & \textbf{0.342} & \textbf{0.326} \\
\hline
\end{tabular}
\caption{Comparison of Reranking models and depths. All experiments used Qdrant, all-minilm-l6-v2 embeddings, and GPT-20b for Query Rewriting.}
\label{tab:rerankers}
\end{table}

As shown in Table \ref{tab:rerankers}, modern rerankers significantly outperform older architectures. While \texttt{MS-Marco} provided negligible gains over the non-reranked pipeline, \texttt{BGE-v2-M3} achieved the highest performance across all metrics. We further investigated the ``re-ranking depth''---the number of documents passed to the reranker. Counter-intuitively, reranking only the top-5 documents yielded higher performance (nDCG@3: 0.342) than reranking the top-10 (nDCG@3: 0.329). This suggests that for this specific benchmark, the documents ranked 6--10 by the vector store act as ``distractors.'' Expanding the window introduces more noise than signal, causing the reranker to falsely promote irrelevant passages.

\subsubsection{Embedding Models and Infrastructure}
We evaluated the semantic capture capabilities of four embedding models. As shown in Table \ref{tab:embeddings}, larger models do not strictly guarantee better performance on the MTRAG dataset. The state-of-the-art \texttt{jina-embeddings-v2} achieved the highest accuracy (0.332), but the compact \texttt{all-minilm} model remained extremely competitive (0.329). Given the negligible performance difference ($<1\%$), \texttt{all-minilm} represents a superior efficiency-performance trade-off for latency-sensitive applications. Conversely, \texttt{snowflake-arctic} suffered a catastrophic drop (0.059), likely due to domain mismatch or tokenization issues.

\begin{table}[h]
\centering
\small
\begin{tabular}{lcc}
\hline
\textbf{Embedding Model} & \textbf{nDCG@3} & \textbf{R@3} \\
\hline
Snowflake-Arctic (137m) & 0.059 & 0.058 \\
Nomic-Embed-Text & 0.227 & 0.219 \\
All-MiniLM-L6-v2 & 0.329 & 0.311 \\
\textbf{Jina-Embeddings-v2} & \textbf{0.332} & \textbf{0.313} \\
\hline
\end{tabular}
\caption{Impact of Embedding Models on retrieval performance. Experiments utilized Qdrant and BGE-v2-M3 (Top-10) reranking.}
\label{tab:embeddings}
\end{table}

Additionally, we performed an infrastructure sanity check by comparing Qdrant against FAISS. Using identical parameters (Nomic embeddings, BGE reranker), FAISS yielded significantly lower accuracy (nDCG@3: 0.050) compared to Qdrant (0.227), highlighting the critical impact of vector store configuration (e.g., quantization, distance metrics) on end-to-end performance.

\subsection{Task B: Generation Performance}

\textbf{Model Comparison:} Here, we present our results of our Task B evaluation inside \ref{tab:task_b_results}. In our experiments, we discovered \texttt{gemma-3-12b} is the best performing model out of 5 models in comparison based on the harmonic mean of  $\mathbf{RL_{F}}$, $\mathbf{RB_{llm}}$ and $\mathbf{RB_{alg}}$ scores.

We can see, \texttt{gpt-oss-20b} is the second best performing model following \texttt{gemma-3-12b}. Even though \texttt{gpt-oss-20b} has a larger size(20B), it is falling behind \texttt{gemma-3-12b} in nearly all metrics, particularly in Faithfulness(0.82 vs 0.74). This may be due to several architectural factors: Gemma-3 utilizes a 5-to-1 interleaved local/global attention mechanism which reduces KV-cache memory pressure compared to standard dense architectures like GPT-oss, allowing the model to maintain higher precision across the long contexts required in multi-turn RAG. Another reason can be the difference between distillation and pre-training quality of models. Gemma probably received extensive instruction tuning on document grounded generation which helps it stay more compliant with referenced documents compared to GPT-oss.  Also, it has been noted in community benchmarks about GPT-oss to have higher hallucination rates. This may be related to its relatively larger size, allowing model to make more creative, fluent, and well-structured responses.
\texttt{ministral-3-14b-reasoning} achieves a balanced performance with high logical coherence, despite its strong reliance on internal reasoning, occasionally compromises strict context adherence with ($\mathbf{RL_{F}}$ =0.76) compared to the top-performing instruction-tuned models.

\texttt{deepseek-r1-0528-qwen3-8b} performs strong algorithmic problem-solving capabilities $\mathbf{RB_{alg}}$ typical of reasoning architectures, this model suffers from reduced faithfulness ($\mathbf{RL_{F}}$=0.75) due to a tendency to prioritize internal chain-of-thought logic over retrieved information.

\textbf{Hallucination Analysis:} As can be seen in the table 1, larger models generally perform better at answering "i do not know" to the unanswerable questions.


\subsection{Task C: End-to-End Performance}
\begin{itemize}
    \item \textbf{Combined Results:} Present the final performance when retrieval and generation are coupled.
    \item \textbf{Error Propagation:} Analyze how retrieval errors in Task A negatively impacted the generation scores in Task C (the "Cascading Error" effect).
\end{itemize}
\begin{table*}[h!]
    \centering
    \small
    \setlength{\tabcolsep}{4pt}
        \begin{tabular}{llcccc} % Updated to 6 columns: l l c c c c
            \hline
            \textbf{Rank} & \textbf{LLM Model} & \textbf{H.Mean} & $\mathbf{RL_{F}}$ & $\mathbf{RB_{llm}}$ & $\mathbf{RB_{alg}}$ \\
            \hline
            1st & Gemma-3-12b & \textbf{0.59} & \textbf{0.82} & 0.69 & \textbf{0.42} \\
            2nd & Gpt-oss-20b & 0.55 & 0.74 & \textbf{0.74} & 0.36 \\
            3rd & Ministral-3-14b-reasoning & 0.54 & 0.76 & 0.68 & 0.36 \\
            4th & Deepseek-r1-0528-qwen3-8b & 0.52 & 0.75 & 0.66 & 0.35 \\
            5th & Ministral-3-3b & 0.51 & 0.79 & 0.71 & 0.31 \\
            \hline
\end{tabular}%
    
    \caption{Task B generation performance including Faithfulness and Reference-Based metrics.}
    \label{tab:task_b_results}
\end{table*}

\subsection{Discussion}
\begin{itemize}
    \item Synthesize the findings to recommend a "Best Model" combination.
    \item Discuss the trade-off: specifically, does using a smaller, quantized embedding model with a strong reranker outperform a heavy embedding model without reranking?
\end{itemize}

\section{Conclusion}
\begin{itemize}
    \item \textbf{Summary:} Recap the proposed pipeline and its performance on the benchmark.
    \item \textbf{Main Contribution:} State clearly that high-performance Multi-Turn RAG requires robust Query Rewriting and optimized two-stage retrieval.
    \item \textbf{Future Improvements:} Mention potential work on adversarial turns or synthetic data augmentation (MTRAG-S).
\end{itemize}

\bibliography{custom}

\appendix

\section{Example Appendix}
\label{sec:appendix}

This is an appendix.


%---------------------------------------------------------
% EXPERIMENTAL RESULTS TABLE
%---------------------------------------------------------
\begin{table*}[ht!]
\centering
\small
\setlength{\tabcolsep}{4pt} % Adjust column spacing to fit width
\begin{tabular}{l l l l c c c c c}
\hline
\textbf{ID} & \textbf{VectorDB} & \textbf{Embedding} & \textbf{Reranker} & \textbf{QR} & \textbf{Top-K} & \textbf{RR-K} & \textbf{nDCG@3} & \textbf{Recall@3} \\
\hline
1 & Faiss & Nomic & BGE-v2-M3 & Yes & 20 & 10 & 0.0503 & 0.0497 \\
2 & Qdrant & MiniLM-L6 & BGE-v2-M3 & Yes & 20 & 5 & 0.3421 & 0.3259 \\
3 & Qdrant & MiniLM-L6 & BGE-v2-M3 & Yes & 20 & 10 & 0.3196 & 0.3032 \\
4 & Qdrant & MiniLM-L6 & MS-MARCO & Yes & 20 & 10 & 0.2932 & 0.2753 \\
5 & Qdrant & MiniLM-L6 & Jina-Rerank-v2 & Yes & 20 & 10 & 0.3162 & 0.2997 \\
6 & Qdrant & MiniLM-L6 & Mxbai-XS & Yes & 20 & 10 & 0.3286 & 0.3114 \\
7 & Qdrant & MiniLM-L6 & BGE-v2-M3 & Yes & 30 & 10 & 0.2272 & 0.2187 \\
8 & Qdrant & Jina-v2-Small & BGE-v2-M3 & Yes & 20 & 10 & 0.3174 & 0.3001 \\
9 & Qdrant & Nomic & BGE-v2-M3 & Yes & 20 & 10 & 0.3317 & 0.3127 \\
10 & Qdrant & Snowflake & BGE-v2-M3 & Yes & 20 & 10 & 0.0587 & 0.0576 \\
11 & Qdrant & MiniLM-L6 & BGE-v2-M3 & No & 20 & 10 & 0.2691 & 0.2553 \\
12 & Qdrant & Nomic & BGE-v2-M3 & No & 20 & 10 & 0.2101 & 0.1977 \\
\hline
\end{tabular}
\caption{Task A Retrieval Performance. Comparison of Vector Databases, Embeddings, and Rerankers. Metrics reported are nDCG@3 and Recall@3. `Top-K` refers to the initial retrieval count, and `RR-K` refers to the count after reranking. QR denotes Query Rewriting.}
\label{tab:task_a_results}
\end{table*}

%%%%%%%%%%%%%%%%%%%%%%%%%

\section{Appendix: Detailed Retrieval Experiments}
\label{sec:appendix_results}

Table \ref{tab:full_results} presents the complete set of experimental configurations and their corresponding performance metrics on the MTRAG dataset. The table is sorted by nDCG@3 to highlight the most effective pipeline combinations.

\begin{table*}[h]
\centering
\scriptsize
\setlength{\tabcolsep}{3pt} % Adjust column spacing to fit
\renewcommand{\arraystretch}{1.2} % Increase row height slightly for readability
\begin{tabular}{|c|c|c|c|c|c|c|c|c|c|c|}
\hline
\textbf{ID} & \textbf{QR} & \textbf{VectorDB} & \textbf{QR Model} & \textbf{Embedding} & \textbf{Reranker} & \textbf{Depth} & \textbf{nDCG@3} & \textbf{nDCG@5} & \textbf{R@3} & \textbf{R@5} \\ \hline
1 & YES & Qdrant & GPT-oss-20b & All-MiniLM-L6 & BGE-v2-M3 & 5 & 0.342 & 0.384 & 0.326 & 0.431 \\
2 & YES & Qdrant & GPT-oss-20b & Jina-Embed-v2 & BGE-v2-M3 & 10 & 0.332 & 0.364 & 0.313 & 0.397 \\
3 & YES & Qdrant & GPT-oss-20b & All-MiniLM-L6 & BGE-v2-M3 & 10 & 0.329 & 0.365 & 0.311 & 0.406 \\
4 & YES & Qdrant & GPT-oss-20b & All-MiniLM-L6 & BGE-v2-M3 & 10 & 0.320 & 0.353 & 0.303 & 0.392 \\
5 & YES & Qdrant & GPT-oss-20b & All-MiniLM-L6 & Mxbai-XS & 10 & 0.317 & 0.351 & 0.300 & 0.389 \\
6 & YES & Qdrant & Gemma-3-12b & All-MiniLM-L6 & BGE-v2-M3 & 5 & 0.317 & 0.359 & 0.302 & 0.405 \\
7 & YES & Qdrant & GPT-oss-20b & All-MiniLM-L6 & Jina-Reranker-v2 & 10 & 0.316 & 0.350 & 0.300 & 0.388 \\
8 & YES & Qdrant & GPT-oss-20b & All-MiniLM-L6 & MS-Marco & 10 & 0.293 & 0.324 & 0.275 & 0.356 \\
9 & YES & Qdrant & GPT-oss-20b & All-MiniLM-L6 & None & 10 & 0.293 & 0.324 & 0.275 & 0.356 \\
10 & YES & Qdrant & Gemma-3-12b & All-MiniLM-L6 & BGE-v2-M3 & 10 & 0.290 & 0.328 & 0.274 & 0.368 \\
11 & YES & Qdrant & Gemma-3-12b & All-MiniLM-L6 & Mxbai-XS & 10 & 0.286 & 0.326 & 0.268 & 0.366 \\
12 & NO & Qdrant & - & All-MiniLM-L6 & BGE-v2-M3 & 10 & 0.269 & 0.300 & 0.255 & 0.334 \\
13 & YES & Qdrant & Gemma-3-12b & All-MiniLM-L6 & None & 10 & 0.268 & 0.295 & 0.253 & 0.324 \\
14 & NO & Qdrant & - & All-MiniLM-L6 & None & 10 & 0.249 & 0.269 & 0.237 & 0.293 \\
15 & YES & Qdrant & GPT-oss-20b & Nomic-Embed & BGE-v2-M3 & 10 & 0.234 & 0.256 & 0.219 & 0.292 \\
16 & YES & Qdrant & Gemma-3-12b & Nomic-Embed & BGE-v2-M3 & 10 & 0.232 & 0.253 & 0.224 & 0.285 \\
17 & NO & Qdrant & - & Nomic-Embed & BGE-v2-M3 & 10 & 0.227 & 0.233 & 0.198 & 0.256 \\
18 & YES & Qdrant & Gemma-3-12b & Nomic-Embed & None & 10 & 0.224 & 0.230 & 0.197 & 0.257 \\
19 & YES & Qdrant & GPT-oss-20b & Nomic-Embed & None & 10 & 0.216 & 0.230 & 0.197 & 0.260 \\
20 & NO & Qdrant & - & Nomic-Embed & None & 10 & 0.211 & 0.208 & 0.174 & 0.227 \\
21 & YES & Qdrant & GPT-oss-20b & Snowflake-137m & BGE-v2-M3 & 10 & 0.059 & 0.067 & 0.058 & 0.078 \\
22 & YES & Faiss & GPT-oss-20b & Nomic-Embed & BGE-v2-M3 & 10 & 0.050 & 0.078 & 0.050 & 0.106 \\
\hline
\end{tabular}
\caption{Full Experimental Results for Task A (Retrieval), sorted by nDCG@3 performance. Abbreviations: QR (Query Rewriting), DB (Database), Depth (Reranking Depth). Rows with missing values in `QR Model` indicate experiments where Query Rewriting was disabled.}
\label{tab:full_results}
\end{table*}


\end{document}

